\begin{table}[H]
\centering
\caption{Nama Daerah Perbandingan Cluster Biplot vs STATIS (Tahunan) (2024)} \makeatletter\@ifundefined{r@tab:app1biplot_statis_comp_name_2024}{\label{tab:app1biplot_statis_comp_name_2024}}{}\makeatother
\begin{adjustbox}{max width=\textwidth}
\begin{tabular}{|p{4cm}|p{10cm}|}
\hline
\textbf{Kategori} & \textbf{Daftar Wilayah} \\
\hline
Biplot 1 (Rendah) \& Statis 1 (Rendah) & Pandeglang, Lebak, Tangerang, Serang, Kota Cilegon, Kota Serang \\
\hline
Biplot 1 (Rendah) \& Statis 2 (Sedang) & - \\
\hline
Biplot 1 (Rendah) \& Statis 3 (Tinggi) & - \\
\hline
Biplot 2 (Sedang) \& Statis 1 (Rendah) & Kep. Seribu, Sukabumi, Cianjur, Garut, Tasikmalaya, Ciamis, Cirebon, Majalengka, Sumedang, Indramayu, Subang, Bandung Barat, Pangandaran, Kota Tasikmalaya, Kota Banjar \\
\hline
Biplot 2 (Sedang) \& Statis 2 (Sedang) & Bogor, Bandung, Kuningan, Purwakarta, Karawang, Kota Cirebon \\
\hline
Biplot 2 (Sedang) \& Statis 3 (Tinggi) & - \\
\hline
Biplot 3 (Tinggi) \& Statis 1 (Rendah) & - \\
\hline
Biplot 3 (Tinggi) \& Statis 2 (Sedang) & Bekasi, Kota Bogor, Kota Sukabumi, Kota Cimahi \\
\hline
Biplot 3 (Tinggi) \& Statis 3 (Tinggi) & Kota Jakarta Selatan, Kota Jakarta Timur, Kota Jakarta Pusat, Kota Jakarta Barat, Kota Jakarta Utara, Kota Bandung, Kota Bekasi, Kota Depok, Kota Tangerang, Kota Tangerang Selatan \\
\hline
\end{tabular}
\end{adjustbox}
\end{table}